%% kapa: exam=1 round=26 year=2015 period=1 form=A subject=economics subject_ko=경제원론 seq=23 number=63
%% source: Q-Net 감정평가사 기출문제 (cst00309) artlSeq=5138001 file=extracted/1008/5138001/제26회_감정평가사 1차_기출/제26회_감정평가사 1차_기출/1교시_감평_A형.pdf
%% conversion: kordoc 4.24.51 -> tools/gen1.py (auto); review=proofread-sonnet; flags=none
\begin{question}{63}
소비이론에 관한 설명으로 옳지 않은 것은?
\begin{choices}
  \item 절대소득가설에 의하면 소비의 이자율탄력성은 0이다.
  \item 절대소득가설에 의하면 기초소비가 있는 경우, 평균소비성향이 한계소비성향보다 크다.
  \item 항상소득가설에 의하면 임시소비는 임시소득에 의해 결정된다.
  \item 상대소득가설에 의하면 장기소비함수는 원점을 통과하는 직선의 형태로 도출된다.
  \item 생애주기가설에 의하면 사람들은 일생에 걸친 소득 변화 양상을 염두에 두고 적절한 소비수준을 결정한다.
\end{choices}
\end{question}
\answer{3}
